%% Essuie-glace : schéma cinématique (quadrilatère articulé O-A-B-D, quatre pivots d'axe z).
%% Chaque tracé est fait dans la couleur de sa classe : c'est ce qui permet au build de regrouper
%% les tracés (g.mech[data-class]) et à l'application de les animer. Les flèches et annotations
%% fixes sont en solideX!60 (teinte différente) et les lettres restent des glyphes (jamais animés).
\documentclass[tikz,border=4pt]{standalone}
\usepackage{liaisons}
\usetikzlibrary{backgrounds}
\begin{document}
\begin{tikzpicture}[x=1cm,y=1cm, solide/.style={line width=1.4pt, line cap=round, line join=round}]
\useasboundingbox (-1.9,-2.15) rectangle (5.6,3.4);  % boîte fixe : sert à l'animation (cm → SVG)
\colorlet{E0}{black}\colorlet{E1}{solideE}\colorlet{E2}{solideA}\colorlet{E3}{solideB}
\coordinate (O) at (0,0);  \coordinate (A) at (0.8,0);
\coordinate (D) at (4,0);  \coordinate (B) at (4,2.2);
\coordinate (T) at (4,3);                            % extrémité du bras (porte-balai)
% Symboles normalisés : quatre pivots d'axe (·, z), vus « en bout »
\pic[s1=E1, s2=E0, liens=false, name=pO] at (O) {pivot bout};   % bâti / manivelle
\pic[s1=E2, s2=E1, liens=false, name=pA] at (A) {pivot bout};   % manivelle / bielle
\pic[s1=E2, s2=E3, liens=false, name=pB] at (B) {pivot bout};   % bielle / balancier
\pic[s1=E3, s2=E0, liens=false, name=pD] at (D) {pivot bout};   % balancier / bâti
\pic at (1.1,-0.9) {bati};  \pic at (2.9,-0.9) {bati};
\begin{scope}[on background layer]
  \draw[solide, E0] (O) -- (0,-0.9) -- (4,-0.9) -- (D);          % 0 : carrosserie + platine
  \draw[solide, E1] (O) -- (A);                                   % 1 : manivelle (motrice)
  \draw[solide, E2] (A) -- (B);                                   % 2 : bielle
  \draw[solide, E3] (D) -- (B) -- (T) (2.95,3) -- (5.05,3);       % 3 : balancier + bras + balai
\end{scope}
% Mouvements (teintes claires : ces tracés ne sont pas animés)
\draw[-{Stealth[length=5pt]}, line width=0.7pt, solideE!60] (190:1.25) arc (190:250:1.25);
\node[solideE!60!black, font=\small] at (216:1.68) {$\omega_{1/0}$};
\draw[{Stealth[length=5pt]}-{Stealth[length=5pt]}, line width=0.7pt, solideB!60]
  ([shift=(86.4:1.15)]D) arc (86.4:130.1:1.15);
\node[solideB!60!black, font=\small] at (4.78,0.62) {$43{,}6^\circ$};
% Étiquettes de classes et points caractéristiques (glyphes : jamais animés)
\node[E0, font=\small\bfseries] at (2.0,-1.3) {E0};
\node[E1, font=\small\bfseries] at (148:1.55) {E1};
\node[E2, font=\small\bfseries] at (2.3,0.55) {E2};
\node[E3, font=\small\bfseries] at (4.75,1.6) {E3};
\node[font=\small] at (-0.42,-0.42) {O};
\node[font=\small] at (1.30,-0.50) {A};
\node[font=\small] at (4.62,2.35) {B};
\node[font=\small] at (4.34,-0.36) {D};
\repere{(4.2,-1.95)}{x}{y}{1}
\end{tikzpicture}
\end{document}
